
\uline{factorisation de} @@M0@@ \uline{comme composé d'une immersion et d'un morphisme lisse}, \uline{le complexe} @@M1@@ \uline{défini en} 2.1 \uline{est strictement parfait}.

Soit @@M2@@, où @@M3@@ est un morphisme lisse, et @@M4@@ une immersion. D'après 1.2, @@M5@@ est une immersion régulière. Si @@M6@@ est l'Idéal définissant @@M7@@, @@M8@@ est donc un @@M9@@-Module localement libre de type fini. De même, comme @@M10@@ est lisse, @@M11@@ est localement libre de type fini, d'où 2.4.

\textbf{Corollaire 2.5.} \uline{Soit} @@M12@@ \uline{un morphisme d'intersection complète lissifiable}, \uline{et} @@M13@@ \uline{une factorisation de} @@M14@@ \uline{par un} @@M15@@-\uline{schéma lisse} @@M16@@. \uline{Soit} @@M17@@ \uline{l'idéal de} @@M18@@ \uline{définissant} @@M19@@. \uline{Alors}, \uline{si} @@M20@@ \uline{désigne le groupe de Grothendieck des} @@M21@@-\uline{Modules localement libres de type fini} (\uline{resp}. \uline{des complexes parfaits sur} @@M22@@), \uline{l'élément} @@M23@@ \uline{défini par} :
@@M24@@
@@M25@@
\uline{ne dépend pas de la factorisation de} @@M26@@ \uline{choisie}. \uline{La valeur en un point} @@M27@@ \uline{de} @@M28@@ \uline{de l'augmentation de} @@M29@@ \uline{n'est autre que la dimension relative virtuelle de} @@M30@@ \uline{sur} @@M31@@ \uline{au point} @@M32@@ (1.9).

L'élément @@M33@@ est bien défini dans les deux cas envisagés d'après 2.4, et sa classe dans @@M34@@ ne dépend pas de la factorisation, d'après 2.2, dans le cas respé, et d'après la démonstration de 2.2 dans le cas non respé. La dernière assertion de 2.5 est triviale.

\textbf{Proposition 2.6.} \uline{Soient} @@M35@@, \uline{et} @@M36@@ \uline{deux morphismes d'intersection complète lissifiables}, \uline{et supposons de plus la condition suivante remplie}:

